Bảng~\ref{tab:leak} cho thấy mức độ của từng dạng rò rỉ khi chỉ đổi đúng một yếu tố so với giao thức chuẩn. Chia theo chu kỳ làm sai số thấp đi 1{,}2--1{,}8 lần dù mô hình và dữ liệu không đổi. Bộ “dự đoán” không dùng đặc trưng nào, chỉ chép nhãn của chu kỳ train gần nhất, có sai số bằng 0{,}01--0{,}15 lần mô hình thật được đánh giá đúng, đúng như Mệnh đề~\ref{md:rowleak}. Chuẩn hoá min--max theo cả vòng đời cell, như trong mã PINN4SOH, làm sai số thấp giả tạo trên XJTU (0{,}80) và HUST (0{,}46), nơi SOH cuối đời của mọi cell gần như bằng nhau nên vị trí trong vòng đời gần như đủ để suy ra SOH; trên TJU (1{,}51) và MIT (2{,}76) nó làm sai số tăng vì xoá mất mức tuyệt đối của đặc trưng. Dù theo chiều nào, con số đo được cũng không phản ánh sai số khi triển khai, vì tuổi thọ của cell đang vận hành chưa biết trước. Lọc ngoại lai dùng cả vòng đời (0{,}87--1{,}01) và chuẩn hoá gồm cả test (0{,}98--1{,}03) tác động nhỏ hơn nhưng vẫn là rò rỉ.

\begin{table}[!htbp]\centering\small
\caption[Kiểm toán rò rỉ]{Kiểm toán rò rỉ: MAE theo cell của cùng một MLP (70\,\% nhãn, 2 lần lặp $\times$ 5 fold) theo giao thức không rò rỉ, và tỉ số giữa MAE dưới mỗi cách làm sai với MAE đó (nhỏ hơn 1: sai số trông thấp hơn thực tế). Cột cuối: chia theo chu kỳ rồi chép nhãn của chu kỳ train gần nhất trong cùng cell, không dùng đặc trưng nào.}\label{tab:leak}
\begin{adjustbox}{max width=\textwidth}\begin{tabular}{lrrrrrr}\toprule
Bộ & MAE ($\times10^{-3}$) & chia theo chu kỳ & min--max theo cell & lọc dùng cả vòng đời & chuẩn hoá gồm cả test & chép nhãn \\\midrule
XJTU & 13{,}11 & 0{,}55 & 0{,}80 & 0{,}87 & 1{,}03 & 0{,}14 \\
TJU & 9{,}86 & 0{,}85 & 1{,}51 & 0{,}98 & 1{,}01 & 0{,}10 \\
MIT & 2{,}80 & 0{,}81 & 2{,}76 & 1{,}01 & 1{,}01 & 0{,}15 \\
HUST & 17{,}46 & 0{,}64 & 0{,}46 & 0{,}96 & 0{,}98 & 0{,}01 \\
\bottomrule\end{tabular}\end{adjustbox}\end{table}
